\section{Tutorial 5: Kernel Tricks}

\subsection{Objectives}

This tutorial develops intuition and technical understanding of kernel methods:

\begin{itemize}
    \item Understand feature maps and implicit representations
    \item Compute and interpret kernel functions
    \item Connect kernels to inner products in feature space
    \item Analyze how kernels induce nonlinear decision boundaries
\end{itemize}

\medskip

\noindent
\textbf{Focus.}  
Understand kernels as a computational and conceptual tool, not just a formula.

\subsection{Exercise 1: Explicit Feature Map}

Consider the polynomial kernel:
\[
k(x, x') = (x^\top x' + 1)^2, \quad x, x' \in \mathbb{R}^2.
\]

\medskip

\noindent
\textbf{Tasks.}
\begin{enumerate}
    \item Expand $k(x,x')$ explicitly.
    \item Find a feature map $\phi(x)$ such that:
    \[
    k(x,x') = \phi(x)^\top \phi(x').
    \]
\end{enumerate}

\medskip

\noindent
\textbf{Solution.}

Let $x = (x_1, x_2)$, $x' = (x_1', x_2')$.

\begin{enumerate}
\item Expand:
\[
(x^\top x' + 1)^2 =
(x_1 x_1' + x_2 x_2' + 1)^2.
\]

\[
= x_1^2 x_1'^2 + x_2^2 x_2'^2 + 1
+ 2 x_1 x_1' x_2 x_2'
+ 2 x_1 x_1' + 2 x_2 x_2'.
\]

\item One valid feature map:
\[
\phi(x) =
\begin{pmatrix}
x_1^2 \\
x_2^2 \\
\sqrt{2} x_1 x_2 \\
\sqrt{2} x_1 \\
\sqrt{2} x_2 \\
1
\end{pmatrix}.
\]
\end{enumerate}

\medskip

\noindent
\textbf{Key takeaway.}  
Kernels correspond to inner products in higher-dimensional spaces.

\subsection{Exercise 2: Kernel Without Explicit Mapping}

Consider the Gaussian kernel:
\[
k(x,x') = \exp(-\gamma \|x - x'\|^2).
\]

\medskip

\noindent
\textbf{Tasks.}
\begin{enumerate}
    \item Explain why computing $\phi(x)$ explicitly is difficult.
    \item Why is the kernel trick useful here?
\end{enumerate}

\medskip

\noindent
\textbf{Solution.}

\begin{enumerate}
\item The feature map $\phi(x)$ is infinite-dimensional.

\item The kernel allows us to compute inner products directly without constructing $\phi(x)$.
\end{enumerate}

\medskip

\noindent
\textbf{Key takeaway.}  
The kernel trick enables computation in high-dimensional (even infinite) spaces efficiently.

\subsection{Exercise 3: Linear vs Kernel Classifier}

Consider two classifiers:

\begin{itemize}
    \item Linear: $f(x) = w^\top x$
    \item Kernel: $f(x) = \sum_i \alpha_i y_i k(x_i, x)$
\end{itemize}

\medskip

\noindent
\textbf{Tasks.}
\begin{enumerate}
    \item What is the key difference between these models?
    \item Why can the kernel classifier represent more complex boundaries?
\end{enumerate}

\medskip

\noindent
\textbf{Solution.}

\begin{enumerate}
\item The kernel classifier depends on training points and pairwise similarities.

\item Because kernels implicitly map data into higher-dimensional spaces, allowing nonlinear separation.
\end{enumerate}

\medskip

\noindent
\textbf{Key takeaway.}  
Kernel methods extend linear models to nonlinear settings.

\subsection{Exercise 4: Decision Boundary Interpretation}

Consider a kernel SVM with Gaussian kernel.

\medskip

\noindent
\textbf{Tasks.}
\begin{enumerate}
    \item Describe qualitatively the decision boundary.
    \item How do support vectors influence the boundary?
\end{enumerate}

\medskip

\noindent
\textbf{Solution.}

\begin{enumerate}
\item The boundary can be highly nonlinear and adapt to local structure.

\item Each support vector contributes a localized influence via the kernel function.
\end{enumerate}

\medskip

\noindent
\textbf{Insight.}  
Kernel methods produce flexible, data-adaptive decision boundaries.

\subsection{Exercise 5: Valid Kernels (Conceptual)}

\textbf{Question.}  
What properties must a function $k(x,x')$ satisfy to be a valid kernel?

\medskip

\noindent
\textbf{Discussion points.}

\begin{itemize}
    \item Symmetry:
    \[
    k(x,x') = k(x',x)
    \]
    \item Positive semi-definiteness:
    \[
    \sum_{i,j} c_i c_j k(x_i, x_j) \geq 0
    \]
\end{itemize}

\medskip

\noindent
\textbf{Insight.}  
Valid kernels correspond to inner products in some feature space.

\subsection{Exercise 6: Kernel Choice and Inductive Bias}

\textbf{Question.}  
How does the choice of kernel affect learning?

\medskip

\noindent
\textbf{Discussion points.}

\begin{itemize}
    \item Polynomial kernel:
    \begin{itemize}
        \item Captures global interactions
    \end{itemize}

    \item Gaussian kernel:
    \begin{itemize}
        \item Captures local similarity
    \end{itemize}
\end{itemize}

\medskip

\noindent
\textbf{Insight.}  
Kernels encode inductive bias through similarity structure.

\subsection{Summary}

\begin{itemize}
    \item Kernels represent inner products in feature space
    \item The kernel trick avoids explicit high-dimensional computation
    \item Kernel methods extend linear models to nonlinear problems
    \item Choice of kernel determines inductive bias and model behavior
\end{itemize}

\medskip

\noindent
\textbf{Preparation for next lecture.}  
We now study how to evaluate and compare models, selecting the best one for a given task.